\documentclass[10pt,a4paper]{amsart}
\usepackage[margin=19mm]{geometry}
\usepackage{amsmath,amssymb,mathtools}
\usepackage[scr=rsfs,cal=euler]{mathalfa}
\usepackage{xfrac}
\usepackage[unicode,hidelinks,bookmarks=false]{hyperref}
\newcommand{\ie}{i.e.\ }
\newcommand{\cf}{cf.\ }
\newcommand{\resp}{respectively\ }
\newcommand{\Subsection}{\subsection}
\providecommand{\nobreakdash}{\nobreak\hbox{-}}
\setlength{\emergencystretch}{2em}
\multlinegap0pt
\usepackage{bm}
\usepackage{bbm}
\usepackage{dsfont}
\usepackage{xspace}
\usepackage{cancel}
\usepackage{mathtools}
\usepackage{amsthm}
\makeatletter
\@ifundefined{lemma}{\newtheorem{lemma}{Lemma}}{}
\makeatother
\def\F{\mathcal{F}}
\def\G{\mathcal{G}}
\def\I{\mathcal{I}}
\def\S{\mathbb{S}}
\def\bN{\mathbb{N}}
\title[RiNNAL+: a Riemannian ALM Solver for SDP-RLT Relaxations of ]{RiNNAL+: a Riemannian ALM Solver for SDP-RLT Relaxations of Mixed-Binary Quadratic Programs}
\author{Di Hou, Tianyun Tang, Kim-Chuan Toh}
\date{Source: arXiv:2507.13776v1 (CC BY 4.0)}
\begin{document}
\maketitle

\noindent\emph{Source and licence.} RiNNAL+: a Riemannian ALM Solver for SDP-RLT Relaxations of Mixed-Binary Quadratic Programs. Version arXiv:2507.13776v1 (CC BY 4.0), distributed under the Creative Commons Attribution 4.0 International licence, as recorded in the arXiv licence metadata for this exact version and shown on the abstract page https://arxiv.org/abs/2507.13776v1. Full rights notice: licenses/arxiv-2507.13776-CC-BY-4.0.txt.\par\smallskip

\noindent\textit{Editorial scope.} Counted unit: the Lemma whose source label is \texttt{lemma-F1-F2} in the article (source0.tex lines 635--637, source label \texttt{lemma-F1-F2}), delivered together with the article's own complete original proof of it (source0.tex lines 638--761, one \texttt{proof} environment of 124 lines). No mathematics is written, repaired or re-derived: statement and proof are the article's own text line for line. Required context retained uncounted: phi (lines 613--631). Every definition, display and label of those lines is retained unchanged. Omitted from this package and not redistributed: 1--612, 632--634, 762--2788. Nothing inside a retained excerpt is omitted or altered: every retained body is delivered exactly as the article's lines are written, so no omission receipt of any kind accompanies this package. Citations: the retained excerpts carry no citation command at all, so no bibliography entry of the article is transcribed and no reference is redistributed. Reference adaptation: none was needed and none was made; every reference inside the delivered unit resolves inside it, so no unresolved reference or citation is produced. Printed slips kept literal: no printed word, symbol, hypothesis or number of the counted statement or of its proof was altered. Layout and class: the article's own class is \texttt{article} with packages including amssymb, mathrsfs, latexsym, amsmath, amsfonts, epsfig, graphicx, cite, psfrag, eepic, color, colordvi, amscd, ebezier; the wrapper is the lane's pinned \texttt{amsart} editorial wrapper at 10pt on A4, which restates the theorem-like environments the retained text needs and adds the article's own macro definitions that the retained text needs (\texttt{\textbackslash F}, \texttt{\textbackslash G}, \texttt{\textbackslash I}, \texttt{\textbackslash S}, \texttt{\textbackslash bN}), with extra wrapper packages (bm, bbm, dsfont, xspace, cancel, mathtools). The delivered numbering is the unit's own. Assets: no retained span contains a figure environment, a graphics-inclusion command, a PDF-page inclusion, a table or a TikZ picture; no image file is copied into this package, the package declares no asset dependency and no image file is redistributed with it. Licence: version arXiv:2507.13776v1 of this article is distributed under the Creative Commons Attribution 4.0 International licence, as recorded in the arXiv licence metadata for this exact version and shown on the abstract page https://arxiv.org/abs/2507.13776v1. A full rights notice is delivered at licenses/arxiv-2507.13776-CC-BY-4.0.txt. Characters: every retained excerpt is pure ASCII, so the delivered engine, XeLaTeX with its default Latin Modern text font, reports no missing character.
     \begin{equation}
     \label{phi}
     \begin{aligned}
     \Phi\left(Y\right):=\begin{bmatrix}
             1&0_{1\times n}\\
             0_{n\times 1}&I_n\\
             d&-G
         \end{bmatrix}Y\begin{bmatrix}
             1&0_{1\times n}\\
             0_{n\times 1}&I_n\\
             d&-G
         \end{bmatrix}^\top
         =\begin{bmatrix}
             z&x^\top&(zd-Gx)^\top\\
             x&X&(dx^\top-GX)^\top\\
             zd-Gx&dx^\top-GX&\G(Y)
         \end{bmatrix},
     \end{aligned}
     \end{equation}

\begin{lemma}\label{lemma-F1-F2}
The mapping $\Phi$ is bijective from ${\mathrm{F}_R}$ to ${\mathrm{F}_D}$.
\end{lemma}

\begin{proof}
It is sufficient to prove that:
\begin{enumerate}
\item[(1)] For any $Y:=\begin{bmatrix}
z&x^\top\\
x&X
\end{bmatrix}\in{\mathrm{F}_R}$, it holds that $\Phi(Y)\in {\mathrm{F}_D}$.
\item[(2)] $\Phi$ is surjective.
\item[(3)] $\Phi$ is injective.
\end{enumerate}
Proof of (1): assume that $Y\in{\mathrm{F}_R}=\F\cap\I_R\cap \mathbb{S}_{+}^{n+1} \cap\mathbb{N}^{n+1}$. Since $z=1$, we have that  
\[
\Phi(Y):=\begin{bmatrix}
1&x^\top&s^\top\\
x&X&Z^\top\\
s&Z&W
\end{bmatrix}=\begin{bmatrix}
             1&x^\top&(d-Gx)^\top\\
             x&X&(dx^\top-GX)^\top\\
             d-Gx&dx^\top-GX&GXG^\top -Gxd^\top -dx^\top G^\top +dd^\top
         \end{bmatrix}.
\]
First, by the definition of $\Phi(Y)$ in \eqref{phi}, $Y\succeq 0$ implies that $\Phi(Y)\succeq 0$. Next, the relation $Y\in\I_R$ implies that $s\geq 0,\ Z\geq 0,\ W\geq 0$. Also, the relation $Y\in \mathbb{N}^{n+1}$ implies that $x\geq 0,\ X\geq 0$. Thus, we have $\Phi(Y)\in \bN^{n+l+1}$. Finally, denote $x':=[x;s]$. 
Then
\begin{equation*}
\begin{aligned}
A'x'
=\begin{bmatrix}
A&0_{m\times l}\\
 G&I_l
\end{bmatrix}\begin{bmatrix}
x\\s
\end{bmatrix}
=\begin{bmatrix}
Ax\\
Gx+s
\end{bmatrix}
=\begin{bmatrix}
b\\
d
\end{bmatrix}
=b', 
\end{aligned}
\end{equation*}
where the third equality follows from the definition of $s$. Similarly,
\[
A'X'
=\begin{bmatrix}
A&0_{m\times l}\\
 G&I_l
\end{bmatrix}\begin{bmatrix}
X&Z^\top\\
Z&W
\end{bmatrix}
=\begin{bmatrix}
AX&AZ^\top\\
GX+Z&GZ^\top+W
\end{bmatrix}
=\begin{bmatrix}
bx^\top&bs^\top\\
dx^\top&ds^\top
\end{bmatrix}
=b'(x')^\top,
\]
where the third equality follows from the definitions of $s$, $Z$ and $W$. Furthermore, 
\[
\operatorname{diag}_B(X')=\operatorname{diag}_B(X)=x_B=x'_B.
\] 
Thus, we have $\Phi(Y)\in \F_D$. According to the relations above, we proved that $\Phi(Y)\in{\mathrm{F}_D}$.
\smallskip

\noindent Proof of (2): we prove that for any $Y'\in{\mathrm{F}_D}$, there exists some $Y\in{\mathrm{F}_R}$ such that $\Phi(Y)=Y'$. Denote
\[
Y':=\begin{bmatrix}
1&x^\top&s^\top\\
x&X&Z^\top\\
s&Z&W
\end{bmatrix}\in {\mathrm{F}_D},
\] 
then it is sufficient to prove that the submatrix $Y_0:=[1, x^\top; x, X] \in{\mathrm{F}_R}$ and it satisfies $\Phi(Y_0)=Y'$. By the relation $Y'\in\F_D$, we have 
\begin{equation}\label{eq-Y'-F'}
\begin{aligned}
\begin{bmatrix}
Ax\\
Gx+s
\end{bmatrix}
=\begin{bmatrix}
b\\
d
\end{bmatrix}
,\quad \begin{bmatrix}
AX&AZ^\top\\
GX+Z&GZ^\top+W
\end{bmatrix}
=\begin{bmatrix}
bx^\top&bs^\top\\
dx^\top&ds^\top
\end{bmatrix},
\end{aligned}
\end{equation}
which implies that
\[
\Phi(Y_0)=\begin{bmatrix}
             1&x^\top&(d-Gx)^\top\\
             x&X&(dx^\top-GX)^\top\\
             d-Gx&dx^\top-GX&GXG^\top -Gxd^\top -dx^\top G^\top +dd^\top
         \end{bmatrix}
=\begin{bmatrix}
1&x^\top&s^\top\\
x&X&Z^\top\\
s&Z&W
\end{bmatrix}=Y'.
\]
Thus, we only need to show that $Y_0\in{\mathrm{F}_R}$.
First, the relation $Y_0\in\S^{n+1}_+\cap\mathbb{N}^{n+1}$ holds because $Y'\in\S^{n+l+1}_+\cap\mathbb{N}^{n+l+1}$. Second, denote $x':=[x;s]$, the relation $Y_0\in\F$ follows from the first $m$ rows of \eqref{eq-Y'-F'} and 
\[
\operatorname{diag}_B(X)=\operatorname{diag}_B(X')=x'_B=x_B.
\]
Finally, the relation $Y_0\in\I_R$ is equivalent to $s\geq0, Z\geq 0, W\geq 0$, which holds because $Y'\geq 0$. Thus, we have $Y_0\in{\mathrm{F}_R}$.
\smallskip

\noindent Proof of (3): we prove that if $\Phi(Y_1)=\Phi(Y_2)$ for any $Y_1,Y_2\in{\mathrm{F}_R}$, then $Y_1=Y_2$.
This can be directly proven by noting that $Y_1$ and $Y_2$ are the submatrices of $\Phi(Y_1)$ and $\Phi(Y_2)$.
\end{proof}

\end{document}
